% Flächeninhalt und Volumen im Raum – bank/flaecheninhalt-und-volumen-im-raum/ (zusammenbau v0.4, Heft)
\einheitenkopf[t10]{Flächeninhalt und Volumen im Raum}
\setcounter{aufgabe}{88}

% Anlauf (ohne Prüfkennung)
\begin{aufgabe}{Dreiecksflächen – Anlauf}
\begin{teile}
\teil Im Dreieck $PQR$ ist $M$ der Mittelpunkt von $PQ$, und $MR$ steht senkrecht auf $PQ$. Was ist $MR$? \\ \kreuz{Höhe zu PQ} \\ \kreuz{schräge Seite}
\teil $P(1 | 1 | 0)$, $Q(4 | 5 | 0)$, $R(-4 | 11 | 0)$, rechtwinklig bei $Q$. Flächeninhalt des Dreiecks $PQR$? A = \leerfeld FE
\teil Die Ebene $3x - 2y + z = 12$ schneidet die $x$-Achse und die $y$-Achse. Diese beiden Schnittpunkte und der Ursprung bilden ein Dreieck. Flächeninhalt? A = \leerfeld FE
\end{teile}
\end{aufgabe}

\begin{aufgabe}{Dreiecksflächen – Prüfungsaufgaben}
\begin{teile}
\teil Der Ursprung $O$, $Q(2 | 4 | 4)$ und $R(5 | 4 | -2)$ bilden ein gleichschenkliges Dreieck mit der Basis $OQ$. Flächeninhalt? \hfill (Abitur 2025 GK) \\ A = \leerfeld FE
\teil Der Ursprung $O$, $P(6 | 8 | 0)$ und $R(3 | 4 | z)$ mit $z > 0$ bilden ein gleichschenkliges Dreieck mit der Basis $OP$ und dem Flächeninhalt 40. Wert von $z$? \hfill (Abitur 2022 GK) \\ z = \leerfeld
\teil $P(1 | 0 | 0)$, $Q(3 | 1 | 0)$, $R(2 | 4 | 1)$; für $T$ gilt $\overrightarrow{OT} = \overrightarrow{OR} - 3 \cdot \overrightarrow{PQ}$. Bestimme $T$ und das Verhältnis der Flächeninhalte von Dreieck $PQR$ und Viereck $PQRT$. \hfill (Abitur 2024 GK)
\end{teile}
\end{aufgabe}

% Anlauf (ohne Prüfkennung)
\begin{aufgabe}{Vierecksflächen – Anlauf}
\begin{teile}
\teil Ein Viereck hat vier gleich lange Seiten, seine Diagonalen sind verschieden lang. Welche Formel für den Flächeninhalt? \\ \kreuz{Seite mal Seite} \\ \kreuz{halbe Summe der parallelen Seiten mal Höhe} \\ \kreuz{halbes Diagonalenprodukt}
\teil $P(1 | 0 | 0)$, $Q(1 | 4 | 3)$, $R(1 | 10 | -5)$. Zeige, dass bei $Q$ ein rechter Winkel liegt, und berechne den Flächeninhalt des Rechtecks $PQRS$. A = \leerfeld FE
\teil Das Viereck mit $P(3 | 4 | 4)$, $Q(-3 | 7 | 4)$, $R(-1 | 2 | 0)$, $S(5 | -1 | 0)$ ist eine Raute. Flächeninhalt? A = \leerfeld FE
\end{teile}
\end{aufgabe}

\begin{aufgabe}{Vierecksflächen – Prüfungsaufgaben}
\begin{teile}
\teil Das Quadrat $PQRS$ mit $P(2 | 1 | 2)$, $Q(8 | 9 | 2)$, $R(0 | 15 | 2)$ und $S(-6 | 7 | 2)$ liegt in der Ebene $z = 2$. Weise nach, dass es den Flächeninhalt 100 hat. \hfill (Abitur 2017 GK)
\teil Das Dach eines Pavillons besteht aus vier gleichen Rauten; eine hat die Ecken $P(6 | 0 | 3)$, $Q(4 | 4 | 5)$, $R(0 | 6 | 3)$, $S(2 | 2 | 1)$ (1 LE = 1 m). Wie groß ist die gesamte Dachfläche? \hfill (Abitur 2022 GK) \\ \leerfeld[m²]
\teil Das Trapez $PQRS$ mit $P(0 | 1 | 1)$, $Q(8 | 1 | 1)$, $R(6 | 4 | 5)$, $S(2 | 4 | 5)$ hat die parallelen Seiten $PQ$ und $SR$ und gleich lange Schenkel. Flächeninhalt? \hfill (Abitur 2019 GK) \\ A = \leerfeld FE
\teil Ein Quadrat liegt in der $x$-$y$-Ebene, eine Ecke ist der Ursprung $O$. Sein Diagonalenschnittpunkt $M$ liegt auf der Geraden $\vec x = (2 | 4 | 4) + t \cdot (2 | 2 | -2)$. Bestimme $M$ und den Flächeninhalt des Quadrats. \hfill (Abitur 2024 GK)
\teil Das Quadrat $PQRS$ mit $P(0 | 1 | 0)$, $Q(8 | 1 | 0)$, $R(8 | 9 | 0)$, $S(0 | 9 | 0)$ wird verändert: $P$ und $S$ werden auf der $y$-Achse gleich weit aufeinander zu verschoben, sodass ein gleichschenkliges Trapez $P'QRS'$ mit drei Vierteln der Quadratfläche entsteht. Bestimme $P'$ und $S'$. \hfill (Abitur 2021 GK)
\teil Ein quadratisches Beet $PQRS$ mit $P(0 | 2 | 0)$, $Q(10 | 2 | 0)$, $R(10 | 12 | 0)$, $S(0 | 12 | 0)$ (1 LE = 1 m) wird verkleinert: $P$ und $S$ rücken auf der $y$-Achse gleich weit aufeinander zu, sodass ein gleichschenkliges Trapez $P'QRS'$ mit drei Fünfteln der Quadratfläche entsteht. Bestimme $P'$ und $S'$. \hfill (Abitur 2021 GK)
\teil Ein Turm hat die quadratische Grundfläche $PQRS$ mit $P(1 | 1 | 0)$, $Q(7 | 1 | 0)$, $R(7 | 7 | 0)$, $S(1 | 7 | 0)$; das Quadrat $KLTU$ mit $K(4 | 1 | 6)$, $L(7 | 4 | 6)$, $T(4 | 7 | 6)$, $U(1 | 4 | 6)$ liegt über den Seitenmitten. Weise nach, dass $PQRS$ doppelt so groß ist wie $KLTU$. \hfill (Abitur 2022 GK)
\teil Ein Sockel hat die quadratische Grundfläche $PQRS$ mit $P(1 | -3 | 0)$, $Q(9 | -3 | 0)$, $R(9 | 5 | 0)$, $S(1 | 5 | 0)$; die Deckfläche $KLTU$ mit $K(5 | -3 | 3)$, $L(9 | 1 | 3)$, $T(5 | 5 | 3)$, $U(1 | 1 | 3)$ liegt über den Seitenmitten. Weise nach, dass $PQRS$ doppelt so groß ist wie $KLTU$. \hfill (Abitur 2022 GK)
\end{teile}
\end{aufgabe}

% Anlauf (ohne Prüfkennung)
\begin{aufgabe}{Ebene Figur: Flächeninhalt des Schattens eines geneigten Rechtecks bei senkrechtem Lichteinfall über eine Querschnittsskizze begründen – Anlauf}
\begin{teile}
\teil Ein rechteckiges Glasdach eines Carports mit der waagerechten Seite 1{,}6\,m und der geneigten Seite 1\,m ist um $40^\circ$ gegen die Waagerechte geneigt; Licht fällt senkrecht auf das Glas. Begründe mit einer Querschnittsskizze, wie lang der Schatten der geneigten Seite auf dem Boden ist, und berechne den Flächeninhalt des Schattens.
\end{teile}
\end{aufgabe}

\begin{aufgabe}{Ebene Figur: Flächeninhalt des Schattens eines geneigten Rechtecks bei senkrechtem Lichteinfall über eine Querschnittsskizze begründen – Prüfungsaufgaben}
\begin{teile}
\teil Ein rechteckiges Vordach ist gegen die Waagerechte geneigt, eine seiner Seiten verläuft waagerecht. Sonnenlicht fällt als parallele Strahlen senkrecht auf die Dachfläche und wirft einen rechteckigen Schatten auf den waagerechten Boden. Begründe mithilfe einer beschrifteten Querschnittsskizze, dass der Schatten einen größeren Flächeninhalt hat als das Vordach. \hfill (Abitur 2017 GK)
\end{teile}
\end{aufgabe}

% Anlauf (ohne Prüfkennung)
\begin{aufgabe}{Volumen und Oberfläche – Anlauf}
\begin{teile}
\teil Ein Zelt über einem quadratischen Boden, alle Dachkanten laufen in einer Spitze zusammen. Volumenformel? \\ \kreuz{mit Faktor ein Drittel} \\ \kreuz{ohne Faktor ein Drittel}
\teil Pyramide mit der Grundfläche $P(1 | 1 | 0)$, $Q(5 | 1 | 0)$, $R(5 | 5 | 0)$, $T(1 | 5 | 0)$ und der Spitze $S(3 | 3 | 6)$. Volumen? V = \leerfeld VE
\teil Eine Pyramide über einem Rechteck mit den Seiten 5 und 8 hat das Volumen 100. Höhe? h = \leerfeld
\end{teile}
\end{aufgabe}

\begin{aufgabe}{Volumen und Oberfläche – Prüfungsaufgaben}
\begin{teile}
\teil Das Dreieck $OPQ$ mit $P(0 | 3 | 4)$ und $Q(6 | 0 | 0)$ ist rechtwinklig in $O$. Die Pyramide $OPQS$ hat das Volumen 40, und es gilt $\overrightarrow{OS} \cdot \overrightarrow{OP} = 0$ und $\overrightarrow{OS} \cdot \overrightarrow{OQ} = 0$. Länge $|OS|$? \hfill (Abitur 2026 GK) \\ |OS| = \leerfeld
\teil Das Quadrat $PQRT$ mit $P(1 | 1 | 2)$, $Q(7 | 1 | 2)$, $R(7 | 7 | 2)$, $T(1 | 7 | 2)$ ist die Grundfläche einer Pyramide $PQRTS$ mit dem Volumen 96. Gib eine mögliche $z$-Koordinate von $S$ an. \hfill (Abitur 2017 GK) \\ z = \leerfeld
\teil Die Pyramide $PQRTS$ hat das Drachenviereck mit $P(0 | 0 | 0)$, $Q(3 | 2 | 0)$, $R(0 | 7 | 0)$, $T(-3 | 2 | 0)$ als Grundfläche und die Spitze $S(0 | 0 | 5)$. Volumen? \hfill (Abitur 2025 GK) \\ V = \leerfeld VE
\teil Ein gerades Prisma hat die Raute mit $P(3 | 0 | 0)$, $Q(0 | 7 | 0)$, $R(-3 | 0 | 0)$, $T(0 | -7 | 0)$ als Grundfläche; die Ecke über $Q$ ist $Q'(0 | 7 | 5)$. Volumen? \hfill (Abitur 2026 GK) \\ V = \leerfeld VE
\teil Ein gerades Prisma hat als Grundfläche ein Dreieck mit den Seiten 5, 12 und 13; seine Mantelfläche ist 150. Höhe des Prismas? \hfill (Abitur 2019 GK) \\ h = \leerfeld
\teil Die Grundfläche der Pyramide $PQRS$ ist ein rechtwinkliges Dreieck $PQR$ mit der Hypotenuse $|PQ| = 13\,\mathrm{cm}$ und der Kathete $|PR| = 5\,\mathrm{cm}$. Die Kante $RS$ steht senkrecht auf der Grundfläche und ist 9\,cm lang. Volumen? \hfill (Abitur 2020 GK) \\ V = \leerfeld[cm³]
\teil Die Grundfläche der Pyramide $PQRS$ ist ein rechtwinkliges Dreieck $PQR$ mit der Hypotenuse $|PQ| = 10\,\mathrm{cm}$ und der Kathete $|PR| = 6\,\mathrm{cm}$. Die Kante $RS$ steht senkrecht auf der Grundfläche und ist 5\,cm lang. Volumen? \hfill (Abitur 2020 GK) \\ V = \leerfeld[cm³]
\end{teile}
\end{aufgabe}

% Anlauf (ohne Prüfkennung)
\begin{aufgabe}{Zusammengesetzt und Verhältnisse – Anlauf}
\begin{teile}
\teil Ein Haus: Quader mit aufgesetztem Satteldach. Wie fasst du sein Volumen? \\ \kreuz{direkt} \\ \kreuz{als Summe} \\ \kreuz{als Differenz}
\teil Ein Quader hat die Grundfläche $PQRU$ mit $P(1 | 1 | 0)$, $Q(7 | 1 | 0)$, $R(7 | 7 | 0)$, $U(1 | 7 | 0)$ und die Höhe 4. Die Ebene durch $Q$, $U$ und $T(1 | 1 | 8)$ zerlegt den Quader. Berechne das Volumen des Teilkörpers, der $P$ enthält, und erläutere den Ansatz. V = \leerfeld VE
\teil Ein Würfel mit der Kantenlänge 8 hat eine Ecke im Ursprung, seine Kanten liegen auf den positiven Achsen. Die Ebene durch die Würfelkante auf der $x$-Achse und den Mittelpunkt $M(8 | 8 | 4)$ einer senkrechten Kante teilt den Würfel. Volumen des kleineren Teilkörpers? V = \leerfeld VE
\end{teile}
\end{aufgabe}

\begin{aufgabe}{Zusammengesetzt und Verhältnisse – Prüfungsaufgaben}
\begin{teile}
\teil Ein Quader hat die Grundfläche $PQRU$ mit $P(2 | 3 | 0)$, $Q(8 | 3 | 0)$, $R(8 | 9 | 0)$, $U(2 | 9 | 0)$ und die Höhe 5. Die Ebene durch $Q$, $U$ und $T(2 | 3 | 10)$ zerlegt den Quader. Berechne das Volumen des Teilkörpers, der $P$ enthält, und erläutere den Ansatz. \hfill (Abitur 2018 GK) \\ V = \leerfeld VE
\teil Ein Sockel hat die Form eines Stumpfs einer geraden quadratischen Pyramide: Grundfläche mit der Seite 6\,m, Deckfläche mit der Seite 3\,m, Höhe 6\,m. Sein Volumen wird mit dem Term $\tfrac{1}{3} \cdot 36 \cdot 12 - \tfrac{1}{3} \cdot 9 \cdot 6$ berechnet. Begründe, dass der Term das Volumen liefert. \hfill (Abitur 2022 GK)
\teil Die Punkte $(9 | 0 | 0)$, $(0 | 9 | 0)$ und $(0 | 0 | 9)$ liegen in der Ebene $x + y + z = 9$; die Punkte $K$, $L$, $N$ liegen je auf einer Koordinatenachse und in der Ebene $3x + 3y + 3z = 12$. Das Volumen des Körpers zwischen den beiden Ebenen im ersten Oktanten wird mit $\tfrac{1}{3} \cdot \tfrac{1}{2} \cdot 9^2 \cdot 9 - \tfrac{1}{3} \cdot \tfrac{1}{2} \cdot k^2 \cdot k$ berechnet. Erläutere den Ansatz und gib $k$ an. \hfill (Abitur 2023 GK)
\teil Ein Quader und eine Pyramide haben dieselbe Grundfläche, die Spitze der Pyramide liegt in der Deckfläche des Quaders. Um wie viel Prozent ist das Volumen des Quaders größer als das der Pyramide? Rechne ohne konkrete Volumina. \hfill (Abitur 2021 GK) \\ \leerfeld \%
\teil Das Dreieck mit den Ecken $(0 | 0)$, $(8 | 0)$ und $(8 | 6)$ rotiert um die $x$-Achse. Oberflächeninhalt des entstehenden Körpers? \hfill (Abitur 2018 GK) \\ O = \leerfeld FE
\teil Eine zylindrische Vorratsdose mit dem Durchmesser 20\,cm fasst 10 Liter. Passt sie in ein Regalfach mit 30\,cm lichter Höhe? \hfill (Abitur 2022 GK)
\end{teile}
\end{aufgabe}

\begin{aufgabe}{Zusammengesetzt und Verhältnisse – Prüfungsaufgaben – weiter}
\begin{teile}
\teil Für $-1 < k < 7$ hat die Pyramide $OP_kQ_kR$ die Ecken $P_k(7 - k | 0 | 0)$, $Q_k(0 | 1 + k | 0)$ und $R(0 | 0 | 3)$. Für welches $k$ ist ihr Volumen am größten? \hfill (Abitur 2023 GK) \\ k = \leerfeld
\teil Für $-4 < k < 8$ hat die Pyramide $OP_kQ_kR$ die Ecken $P_k(4 + k | 0 | 0)$, $Q_k(0 | 8 - k | 0)$ und $R(0 | 0 | 9)$. Für welches $k$ ist ihr Volumen am größten? \hfill (Abitur 2023 GK) \\ k = \leerfeld
\teil Ein Quader hat eine quadratische Grundfläche mit der Seite 6 in der $x$-$y$-Ebene und die Höhe 8. Auf der senkrechten Achse durch seine Mitte liegen $S_k$ in der Höhe $4 - k$ und $Q_k$ in der Höhe $4 + k$ ($0 \le k < 4$). Die Pyramide über der Grundfläche mit der Spitze $S_k$ und die Pyramide über der Deckfläche mit der Spitze $Q_k$ haben zusammen 25\,\% des Quadervolumens. Wert von $k$? \hfill (Abitur 2026 GK) \\ k = \leerfeld
\teil Ein Quader hat eine rechteckige Grundfläche mit den Seiten 4 und 5 in der $x$-$y$-Ebene und die Höhe 10. Auf der senkrechten Achse durch seine Mitte liegen $S_k$ in der Höhe $5 - k$ und $Q_k$ in der Höhe $5 + k$ ($0 \le k < 5$). Die Pyramide über der Grundfläche mit der Spitze $S_k$ und die Pyramide über der Deckfläche mit der Spitze $Q_k$ haben zusammen 20\,\% des Quadervolumens. Wert von $k$? \hfill (Abitur 2026 GK) \\ k = \leerfeld
\teil Aus einem Quader mit der Grundfläche 12 und der Höhe 6 wird die Pyramide über der Grundfläche mit der Spitze $S_t$ in der Höhe $6 - 2t$ herausgenommen ($0 \le t \le 3$). Die Abbildung zeigt zwei Graphen, einer gibt das Volumen des Restkörpers in Abhängigkeit von $t$ an. Welcher? Begründe. \hfill (Abitur 2023 GK) \\ \kreuz{Graph I} \\ \kreuz{Graph II}

\begin{ksys}[xmin=0,xmax=4,ymin=0,ymax=80,xstep=1,ystep=10,xlabel=t,ylabel=V] \funktionab{48+8*\x}{II}{0}{3} \funktionab{48+8/3*\x^2}{I}{0}{3} \end{ksys}
\teil Aus einem Quader mit der Grundfläche 12 und der Höhe 9 wird die Pyramide über der Grundfläche mit der Spitze $S_t$ in der Höhe $9 - 3t$ herausgenommen ($0 \le t \le 3$). Die Abbildung zeigt zwei Graphen, einer gibt das Volumen des Restkörpers in Abhängigkeit von $t$ an. Welcher? Begründe. \hfill (Abitur 2023 GK) \\ \kreuz{Graph I} \\ \kreuz{Graph II}

\begin{ksys}[xmin=0,xmax=4,ymin=0,ymax=120,xstep=1,ystep=20,xlabel=t,ylabel=V] \funktionab{72+12*\x}{I}{0}{3} \funktionab{72+4*\x^2}{II}{0}{3} \end{ksys}
\end{teile}
\end{aufgabe}
